\section{Related Work}
\label{sec:related_work}

% Target: ~0.75 two-column page (~1.5 single-column equivalent)
% Condense report sections/related_work/*.tex and presentation slides.
% One subsection per approach, concluding with the research gap.

\subsection{User-Level Mitigations}
\label{sub:user_level}

% [~0.25 column]
%   - FaCT (Cauligi et al. 2019): DSL that guarantees constant-time
%     execution w.r.t. secret-labeled data. Compiler enforces data-flow
%     restrictions.
%   - Limitations: cannot control hardware variability (execution-unit
%     contention, cache interference), compiler optimisations can
%     inadvertently introduce branches, and the OS scheduler itself
%     remains an uncontrolled channel.
%   - Conclusion: user-level is necessary but insufficient for full time
%     protection.

\subsection{seL4 Time Protection}
\label{sub:sel4}

% [~0.5 column]
%   - Ge et al. (2019): extended seL4 with time protection by (i) cloning
%     kernel code into each domain, (ii) flushing shared on-core state that
%     cannot be spatially partitioned, and (iii) using cache colouring for
%     higher-level caches too large to flush.
%   - Five requirements: flush on-core state, partition the OS, enforce
%     deterministic data sharing, flush deterministically, partition
%     interrupts.
%   - Threat model: Trojan in a separate domain from spy; covert and side
%     channels via shared microarchitectural state.
%   - Key difference from this work: seL4 is designed from scratch with
%     capabilities and formal verification; this work retrofits into an
%     existing mainstream kernel.

\subsection{S3K}
\label{sub:s3k}

% [~0.5 column]
%   - Karlsson et al. (2025): capability-based multicore partitioning
%     kernel for RISC-V with time protection as a first-class primitive.
%   - Time-driven scheduler: major/minor frame scheme with time-slice
%     capabilities that can be subdivided and delegated.
%   - Entire kernel resides in scratchpad memory (SPM) backed by L1 cache
%     to avoid leaving cache footprints. fence.t used on every context
%     switch. Constant-time system calls by design.
%   - Evaluation demonstrates no measurable information flow between Trojan
%     and spy when SPM is used.
%   - Key difference from this work: S3K is a new kernel, not a retrofit
%     into an existing OS. FreeRTOS is orders of magnitude more widely
%     deployed.

\subsection{Research Gap}
\label{sub:research_gap}

% [~0.15 column -- transition paragraph]
%   - Both seL4 and S3K demonstrate that time protection is architecturally
%     feasible when designed from the ground up. Neither addresses the
%     fundamentally different challenge of integrating these mechanisms
%     into an existing, widely deployed OS.
%   - This paper fills that gap by showing how domain scheduling and
%     temporal flushing can be added to FreeRTOS without a complete
%     redesign.
